\chapter{Estatística aplicada a retornos}
\label{cap:estatistica}

\abertura{Dez anos de retornos diários de uma ação são cerca de 2\,500 números. Ninguém decide nada olhando para 2\,500 números; decide-se olhando para meia dúzia de resumos deles e para o resultado de alguns testes. Este capítulo trata de quais resumos calcular, de qual distribuição esperar e de quanto se pode confiar em cada conclusão, porque em finanças a amostra quase sempre é menor do que parece.}

\begin{objetivos}
  \item resumir uma série de retornos com medidas de posição, dispersão, assimetria e curtose, sabendo o que cada biblioteca devolve;
  \item explicar por que a distribuição normal subestima a frequência de retornos extremos e usar a $t$ de Student como alternativa;
  \item ler um histograma em escala logarítmica e um QQ-plot;
  \item interpretar corretamente um p-valor e reconhecer quando um teste tem pouco poder;
  \item aplicar os testes $t$, de Jarque-Bera e de Ljung-Box a retornos;
  \item calcular covariâncias e correlações, e desconfiar da estabilidade delas.
\end{objetivos}

\section{Resumir uma série de retornos}
\label{sec:c03-descritiva}

O Código~\ref{cod:c03-descritiva} calcula, para PETR4, VALE3, ITUB4, WEGE3 e BOVA11 entre 2015 e 2024, os resumos que você vai pedir primeiro a qualquer série nova. Os log-retornos estão multiplicados por 100, isto é, em porcentagem ao dia.

\lstinputlisting[style=py, caption={Medidas de posição e de dispersão (\arq{cap03/l01\_descritiva.py})}, label=cod:c03-descritiva]{scripts/cap03/l01_descritiva.py}
\begin{lstlisting}[style=saida]
2461 retornos diários por ativo, em %
        média  mediana  desvio   IQR  IQR/1,349   q5%  q95%    mín    máx
PETR4    0.10     0.13    3.03  2.86       2.12 -4.40  4.50 -35.24  20.07
VALE3    0.06     0.00    2.77  2.78       2.06 -4.03  4.45 -28.18  19.36
ITUB4    0.05     0.03    1.92  2.14       1.59 -2.91  3.10 -12.84  10.49
WEGE3    0.10     0.03    2.12  2.24       1.66 -3.11  3.42 -23.09  13.01
BOVA11   0.04     0.08    1.55  1.62       1.20 -2.28  2.25 -15.75  12.57
\end{lstlisting}

As duas primeiras colunas são \textbf{medidas de posição}\index{medida de posição}, que localizam o centro da distribuição. A \termo{média} amostral $\bar r = \frac{1}{n}\sum_t r_t$ usa todos os valores e por isso é sensível a dias extremos; a \termo{mediana}, o valor que deixa metade das observações de cada lado, não é. As duas são minúsculas, de no máximo $0,13\%$ ao dia, e a mediana de VALE3 é exatamente zero porque em 15 pregões o fechamento da ação não mudou, e esses zeros caem no meio da distribuição. O \cod{dropna()} descarta as 16 datas em que falta o preço de algum dos cinco papéis, e o retorno seguinte a cada uma cobre mais de um pregão; nenhum mínimo ou máximo da tabela vem de um desses casos.

As colunas seguintes medem \termo{dispersão}. O \termo{desvio-padrão} amostral,
\begin{equation}
  s = \sqrt{\frac{1}{n-1}\sum_{t=1}^{n} (r_t - \bar r)^2} ,
  \label{eq:c03-desvio}
\end{equation}
é a volatilidade\index{volatilidade} diária do Capítulo~\ref{cap:retornos}. Compare-o com a média: em PETR4, $0,10\%$ contra $3,03\%$. O retorno esperado de um dia é cerca de trinta vezes menor do que a oscilação típica desse dia, e essa desproporção é a origem de quase toda a dificuldade estatística em finanças. Ela volta na Seção~\ref{sec:c03-media}.

Um \termo{quantil} de ordem $p$ é o valor abaixo do qual fica uma fração $p$ dos dados. A coluna \cod{q5\%} diz que, em um a cada vinte pregões, o log-retorno de PETR4 ficou abaixo de $-4,40\%$. Esse número, com o sinal trocado, é a ideia por trás do \ingles{Value at Risk}, que o Capítulo~\ref{cap:risco} desenvolve. A distância entre os quantis de 75\% e de 25\% é o \termo{intervalo interquartil} (IQR), uma medida de dispersão que ignora o que acontece nas caudas.

Em uma distribuição normal, o IQR vale $1,349$ desvio-padrão. Dividir o IQR por $1,349$ dá, portanto, uma segunda estimativa de $\sigma$, que deveria coincidir com $s$ se os dados fossem normais. Não coincide em nenhuma linha: para PETR4 ela é $2,12\%$, contra $s = 3,03\%$. O miolo da distribuição é mais estreito do que o desvio-padrão sugere, e o desvio-padrão é grande porque uns poucos dias muito ruins (o mínimo de $-35,24\%$ é o 9 de março de 2020 da Figura~\ref{fig:preco-retorno}, uma queda de $29,7\%$ em retorno simples) pesam muito em uma soma de quadrados.

\section{Assimetria e curtose}
\label{sec:c03-momentos}

Média e desvio-padrão descrevem onde a distribuição está e quão larga ela é. Para descrever a forma, usam-se os dois momentos seguintes da variável padronizada $(r - \mu)/\sigma$.

\begin{definicao}{Assimetria e curtose}
Para uma variável aleatória $X$ com média $\mu$ e desvio-padrão $\sigma$,
\begin{equation}
  S = \E\!\left[\left(\frac{X-\mu}{\sigma}\right)^{\!3}\right] , \qquad
  K = \E\!\left[\left(\frac{X-\mu}{\sigma}\right)^{\!4}\right] .
  \label{eq:c03-momentos}
\end{equation}
$S$ é a \termo{assimetria} e $K$ é a \termo{curtose}. Toda distribuição normal tem $S = 0$ e $K = 3$. A diferença $K - 3$ é o \termo{excesso de curtose}.
\end{definicao}

O cubo preserva o sinal do desvio, de modo que $S < 0$ indica que os desvios grandes ocorrem mais para baixo do que para cima: a cauda esquerda é mais longa. A quarta potência elimina o sinal e amplifica ainda mais os desvios grandes. Um dia a 5 desvios-padrão da média contribui com $5^4 = 625$ para a soma, o mesmo que 625 dias a 1 desvio. A curtose mede, por isso, o peso das caudas, e uma distribuição com $K > 3$ é dita \textbf{leptocúrtica}\index{leptocurtose} ou de \termo{caudas pesadas}.

O Código~\ref{cod:c03-momentos} calcula os dois momentos para o Ibovespa. Daqui em diante as listagens omitem as importações e a leitura do arquivo, iguais às do Código~\ref{cod:c03-descritiva}; os arquivos completos estão em \arq{livro/scripts/cap03/}.

\lstinputlisting[style=py, caption={Assimetria e curtose do Ibovespa, de 2010 a 2025 (\arq{cap03/l02\_momentos.py})}, label=cod:c03-momentos, firstline=7]{scripts/cap03/l02_momentos.py}
\Needspace{7\baselineskip}
\begin{lstlisting}[style=saida]
n = 3963
assimetria  pandas -0.7802   scipy -0.7799
curtose     pandas 12.0940   scipy 12.0772
            scipy com fisher=False 15.0772
semanal  n =  835   assimetria -0.51   excesso de curtose  5.22
mensal   n =  192   assimetria -0.97   excesso de curtose  5.32
\end{lstlisting}

\begin{armadilha}{Curtose ou excesso de curtose?}
Tanto \cod{Series.kurt()} do pandas quanto \cod{scipy.stats.kurtosis} devolvem, por padrão, o \emph{excesso} de curtose (a normal dá zero). Para obter $K$ é preciso pedir \cod{fisher=False} no SciPy ou somar 3. As duas bibliotecas também diferem ligeiramente porque o pandas aplica uma correção de viés para amostras finitas e o SciPy, por padrão, não. Ao ler ``curtose de 12'' em um relatório, descubra qual das duas definições foi usada.
\end{armadilha}

O Ibovespa diário tem assimetria de $-0,78$ e excesso de curtose de $12,1$. As duas últimas linhas da saída repetem a conta com retornos semanais e mensais, obtidos somando os log-retornos diários. O excesso de curtose cai para perto de 5, mas não desaparece.

Antes de tirar conclusões desses números, veja o quanto eles dependem de poucos dias. Sem os 248 pregões de 2020, o excesso de curtose diário cai de $12,1$ para $2,0$ e a assimetria, de $-0,78$ para $-0,18$. Na série mensal, basta retirar um único mês, março de 2020, com log-retorno de $-35,5\%$, para o excesso de curtose ir de $5,3$ a $-0,2$.\nota{Esses números, e os demais ``sem 2020'' do capítulo, saem de \arq{cap03/\allowbreak c03\_sensibilidade.py}.} Estimativas de momentos de ordem alta são dominadas pelas maiores observações da amostra, e seu erro de estimação é grande. Isso não autoriza apagar a crise dos dados, já que ela é exatamente o tipo de evento que a gestão de risco existe para enfrentar. Autoriza desconfiar de qualquer comparação fina de curtoses entre ativos ou entre períodos.

\section{A normal como ponto de partida, e onde ela falha}
\label{sec:c03-normal}

A distribuição normal\index{distribuicao normal@distribuição normal} $N(\mu, \sigma^2)$ tem densidade
\begin{equation}
  f(x) = \frac{1}{\sigma\sqrt{2\pi}} \exp\!\left[-\frac{(x-\mu)^2}{2\sigma^2}\right]
  \label{eq:c03-normal}
\end{equation}
e é o modelo de referência por boas razões. Dois parâmetros a descrevem por completo; somas de normais independentes são normais, o que combina com a aditividade dos log-retornos; e o teorema central do limite diz que somas de muitas parcelas independentes de variância finita tendem a ela. Se o retorno de um dia é a soma de milhares de pequenas variações, a normal parece uma hipótese natural.

O problema está no expoente. A densidade cai com o \emph{quadrado} da distância à média, o que torna eventos distantes praticamente impossíveis: $95,45\%$ da probabilidade fica a menos de 2 desvios, $99,73\%$ a menos de 3, e um desvio de mais de 5 tem probabilidade de $5,7 \times 10^{-7}$. A Figura~\ref{fig:c03-hist} confronta essa previsão com os $3\,963$ retornos diários do Ibovespa.

\figura[t]{cap03/hist_log}{Log-retornos diários do Ibovespa, de 2010 a 2025, e a normal com mesma média e mesmo desvio-padrão. Em escala linear o ajuste parece razoável, com um pico mais alto e mais estreito. Em escala logarítmica a densidade normal vira uma parábola, e ficam visíveis os dias além de $\pm 6\%$, para os quais ela atribui probabilidade praticamente nula.}{fig:c03-hist}

O histograma em escala linear esconde justamente a região que importa para o risco, porque um dia isolado em três mil não tem altura visível. Trocar o eixo vertical para escala logarítmica é o primeiro diagnóstico de cauda: o logaritmo da densidade normal é uma parábola, e tudo o que aparece fora dela é excesso em relação à normal. O Código~\ref{cod:c03-caudas}, na próxima seção, conta esses dias. Além de 4 desvios-padrão a normal previa $0,25$ ocorrência em toda a amostra; houve 15.

Esse comportamento não é particular do Ibovespa nem de 2020. \citet{mandelbrot1963} já o documentava em preços de algodão, e \citet{cont2001} reuniu sob o nome de \termo{fatos estilizados} as regularidades que aparecem em praticamente todos os ativos e mercados. Quatro delas organizam o resto deste capítulo e o próximo.

\begin{itemize}
  \item \textbf{Caudas pesadas.} A distribuição incondicional dos retornos tem mais massa nos extremos do que a normal. Cont relata índices de cauda entre 2 e 5 na maioria das séries que estudou, o que é compatível com variância finita e deixa em dúvida a existência do quarto momento.
  \item \textbf{Ausência de autocorrelação linear.} O retorno de hoje praticamente não ajuda a prever, por uma regra linear, o de amanhã, exceto em escalas intradiárias muito curtas.
  \item \textbf{Agrupamento de volatilidade.} Dias de grande oscilação, para cima ou para baixo, tendem a vir juntos. O \emph{tamanho} dos retornos é autocorrelacionado mesmo quando os retornos não são.
  \item \textbf{Normalidade agregacional.} Quanto maior o intervalo sobre o qual se calcula o retorno (semana, mês, ano), mais a distribuição se aproxima da normal.
\end{itemize}

A Figura~\ref{fig:c03-acf} mostra o segundo e o terceiro fatos. A \termo{autocorrelação} de ordem $k$ é a correlação da série com ela mesma defasada de $k$ dias. Para os retornos, ela nunca passa de $0,08$ em módulo, embora saia em algumas defasagens da faixa em que estaria por puro acaso (a Seção~\ref{sec:c03-jb-lb} volta a esse ponto); para os retornos ao quadrado, que medem o tamanho do movimento sem o sinal, ela começa em $0,5$ e leva semanas para se dissipar.

\figura[t]{cap03/acf}{Autocorrelação dos log-retornos diários do Ibovespa e dos seus quadrados, de 2010 a 2025. O sinal do retorno de amanhã é quase imprevisível a partir do passado; o tamanho, não. A faixa cinza é $\pm 1,96/\sqrt{n}$.}{fig:c03-acf}

O agrupamento explica por que a normalidade agregacional é lenta nos nossos dados. A versão mais simples do teorema central do limite\index{teorema central do limite} pede parcelas independentes. Com a dependência que os retornos têm (não correlacionados, mas com volatilidade agrupada e variância finita) ele continua valendo, só que a convergência é mais lenta: dos nove dias da amostra com mais de 5 desvios-padrão, sete ocorreram no mesmo mês, março de 2020. Dias extremos que chegam em bloco não se compensam dentro do mês. O Capítulo~\ref{cap:series} modela essa dependência com a família GARCH.

\section{A distribuição $t$ de Student}
\label{sec:c03-t}

Uma distribuição com caudas pesadas, um único parâmetro a mais do que a normal e disponível em qualquer biblioteca é a $t$ de Student.

\begin{definicao}{Distribuição $t$ de Student}
Uma variável $T$ tem distribuição $t$ com $\nu > 0$ \termo{graus de liberdade} se sua densidade é
\begin{equation}
  f(t) = c_\nu \left(1 + \frac{t^2}{\nu}\right)^{-\frac{\nu+1}{2}} , \qquad
  c_\nu = \frac{\Gamma\!\left(\frac{\nu+1}{2}\right)}{\sqrt{\nu\pi}\;\Gamma\!\left(\frac{\nu}{2}\right)} .
  \label{eq:c03-t}
\end{equation}
A média é zero para $\nu > 1$. A variância é $\nu/(\nu-2)$ e só existe para $\nu > 2$. O excesso de curtose é $6/(\nu-4)$ e só existe para $\nu > 4$.
\end{definicao}

A diferença em relação à Equação~\eqref{eq:c03-normal} é a forma da queda. Para $|t|$ grande, a densidade da $t$ se comporta como $|t|^{-(\nu+1)}$, uma potência, enquanto a da normal cai como a exponencial de um quadrado. Uma potência decresce muito mais devagar, e quanto menor $\nu$, mais devagar. É por isso que os momentos deixam de existir: a integral que define o momento de ordem $k$ só converge se $k < \nu$. No outro extremo, quando $\nu \to \infty$ a $t$ converge para a normal padrão, e com $\nu = 30$ as duas já são difíceis de distinguir a olho.

Para comparar caudas de forma justa é preciso igualar as variâncias. Como $\Var(T) = \nu/(\nu-2)$, a variável $X = T\sqrt{(\nu-2)/\nu}$ tem variância 1; é a \termo{$t$ padronizada}. A Figura~\ref{fig:c03-t} e a Tabela~\ref{tab:c03-caudas} comparam a normal padrão com a $t$ padronizada.

\figura[t]{cap03/t_student}{Normal padrão e $t$ de Student padronizada para variância 1. À esquerda, a $t$ concentra mais massa no centro e nas caudas, e menos nos ``ombros''. À direita, em escala logarítmica, a probabilidade de um valor acima de $k$ desvios-padrão: as curvas se cruzam perto de $k = 2$ e depois se separam por ordens de grandeza.}{fig:c03-t}

\begin{table}[h]
\caption{Uma queda de mais de $k$ desvios-padrão ocorre, em média, uma vez a cada quantos pregões? Valores de $1/P(X < -k)$ para distribuições com variância 1.}
\label{tab:c03-caudas}
\begin{tabularx}{\textwidth}{l >{\raggedleft\arraybackslash}X r r r}
\toprule
\cab{$k$} & \cab{Normal} & \cab{$t$, $\nu=10$} & \cab{$t$, $\nu=5$} & \cab{$t$, $\nu=3$} \\
\midrule
2 & 44 & 41 & 41 & 49 \\
3 & 741 & 273 & 171 & 144 \\
4 & 31\,574 & 1\,676 & 560 & 324 \\
5 & 3\,488\,556 & 8\,665 & 1\,506 & 617 \\
6 & 1\,013\,594\,692 & 37\,663 & 3\,489 & 1\,052 \\
\bottomrule
\end{tabularx}
\end{table}

Em 2 desvios-padrão as quatro distribuições dão probabilidades parecidas, e a $t$ com $\nu = 3$ dá até um pouco menos do que a normal. A diferença aparece depois. Sob a normal, uma queda de 5 desvios acontece uma vez a cada $3,5$ milhões de pregões, quase 14 mil anos de bolsa. Sob uma $t$ com $\nu = 5$, uma vez a cada $1\,506$ pregões, ou seis anos. Quem dimensiona risco com a normal trata como impossível algo que um investidor vê algumas vezes na vida.\nota{A tabela e os números deste parágrafo saem de \arq{cap03/\allowbreak c03\_tabela\_caudas.py}.}

\subsection{Ajuste aos dados}

Retornos não têm média zero nem variância $\nu/(\nu-2)$, de modo que na prática se usa a família de \termo{locação e escala} $r = \mu + \lambda\,T$. São três parâmetros: $\nu$, a locação $\mu$ e a escala $\lambda$. O Código~\ref{cod:c03-caudas} os estima por máxima verossimilhança e compara, para cada $k$, o número de dias a mais de $k$ desvios-padrão da média com o que cada distribuição previa.

\lstinputlisting[style=py, caption={Ajuste da $t$ de Student e contagem de dias extremos (\arq{cap03/l03\_caudas.py})}, label=cod:c03-caudas, firstline=7]{scripts/cap03/l03_caudas.py}
\begin{lstlisting}[style=saida]
amostra:    n = 3963  média = 0.00021  desvio = 0.01463
t ajustada: nu = 4.86  loc = 0.00040  escala = 0.01089  desvio = 0.01420
 k   observado   normal   t ajustada
 2         157   180.32       177.68
 3          39    10.70        42.21
 4          15     0.25        13.03
 5           9     0.00         4.93
\end{lstlisting}

O ajuste devolve $\hat\nu = 4,86$, dentro da faixa relatada por \citet{cont2001}. Repare que a escala, $0,01089$, não é o desvio-padrão: pela fórmula da variância, o desvio-padrão da $t$ ajustada é $\lambda\sqrt{\nu/(\nu-2)} = 0,0142$, próximo do amostral. Confundir os dois parâmetros ao simular ou ao calcular quantis subestima a volatilidade em cerca de um quarto.

A $t$ acerta a contagem de dias além de 3 e de 4 desvios, que a normal erra por fatores de $3,6$ e de 60. Além de 5 desvios ela ainda prevê menos (cerca de 5) do que se observou (9). Parte dessa diferença vem do agrupamento: uma distribuição única para todos os dias, por mais pesada que seja a cauda, não descreve o fato de que sete daqueles nove dias foram vizinhos.

\subsection{QQ-plot}

O \termo{QQ-plot} (gráfico quantil-quantil) compara uma amostra com uma distribuição candidata ponto a ponto. Ordenam-se as $n$ observações; a $i$-ésima menor é o quantil empírico de ordem aproximadamente $(i - 0,5)/n$; no eixo horizontal vai o quantil de mesma ordem da distribuição teórica. Se a candidata descreve bem os dados, os pontos caem sobre a reta de $45^\circ$. Caudas pesadas produzem um S: os pontos ficam abaixo da reta à esquerda e acima à direita, porque os extremos observados são maiores do que os teóricos.

A Figura~\ref{fig:c03-qq} resume a seção. O QQ-plot tem uma vantagem sobre o histograma: não depende da escolha de classes, e cada observação extrema aparece individualmente.

\figura[b]{cap03/qq}{QQ-plots dos log-retornos diários do Ibovespa. Contra a normal, os pontos formam o S típico de caudas pesadas, e a pior queda ($-16,0\%$) aparece onde a normal esperava $-5,3\%$. Contra a $t$ ajustada, o alinhamento vai até perto de $\pm 6\%$; os dias de março de 2020 continuam fora.}{fig:c03-qq}

\begin{repositorio}
\arq{aula3/s3\_01\_estatistica\_descritiva.py} monta a tabela descritiva para PETR4, VALE3 e ITUB4, e \arq{aula3/s3\_02\_distribuicoes.py} ajusta a normal e a $t$ ao Ibovespa e compara os quantis de cada uma. Os dois baixam os dados do Yahoo Finance na hora, e por isso os números variam um pouco em relação aos desta página.
\end{repositorio}

\section{A lógica de um teste de hipótese}
\label{sec:c03-teste}

Um teste de hipótese transforma a impressão deixada por um gráfico em uma probabilidade. O raciocínio é o mesmo em todos os testes deste capítulo. Formula-se uma \termo{hipótese nula} $H_0$, uma afirmação precisa sobre o processo que gerou os dados (``o retorno esperado é zero'', ``a distribuição é normal''), e uma alternativa $H_1$. Escolhe-se uma \termo{estatística de teste}, um número calculado da amostra cuja distribuição é conhecida \emph{se $H_0$ for verdadeira}. Calcula-se a estatística nos dados e pergunta-se quão surpreendente é o valor obtido.

\begin{definicao}{p-valor}
O \termo{p-valor} é a probabilidade, calculada supondo $H_0$ verdadeira, de a estatística de teste assumir um valor tão ou mais extremo do que o observado.
\end{definicao}

Um p-valor pequeno diz que os dados seriam raros em um mundo no qual $H_0$ vale, e por isso depõe contra ela. Fixa-se antes um \termo{nível de significância} $\alpha$, por convenção 5\%, e rejeita-se $H_0$ quando o p-valor fica abaixo dele.

\begin{armadilha}{O que o p-valor não é}
O p-valor não é a probabilidade de $H_0$ ser verdadeira, nem a probabilidade de o resultado ter ocorrido ``por acaso''. Ele é calculado \emph{dentro} da hipótese nula e nada diz sobre a chance dela. Também não mede o tamanho do efeito: com amostra grande, uma diferença irrelevante na prática produz p-valor minúsculo. E um p-valor alto não prova $H_0$; pode apenas indicar que a amostra é pequena demais para enxergar o efeito.
\end{armadilha}

A decisão pode errar de duas maneiras. O \termo{erro tipo I} é rejeitar uma $H_0$ verdadeira, um alarme falso, e sua probabilidade é $\alpha$ por construção. O \termo{erro tipo II} é não rejeitar uma $H_0$ falsa, com probabilidade $\beta$. O \termo{poder} do teste é $1 - \beta$, a probabilidade de detectar um efeito que existe. O painel esquerdo da Figura~\ref{fig:c03-poder} mostra as três áreas. O poder cresce com o tamanho do efeito e com o tamanho da amostra, e cai com a dispersão dos dados. Em finanças os efeitos são pequenos e a dispersão é grande, e por isso o poder costuma ser baixo.

\figura[t]{cap03/poder}{À esquerda, a distribuição da estatística $t$ sob $H_0$ e sob uma alternativa em que seu valor esperado é 2; as linhas tracejadas marcam $\pm 1,96$. À direita, o poder do teste de média zero em função dos anos de dados, para três valores da razão entre retorno e volatilidade anuais. Com razão de $0,5$, são precisos 31 anos para um poder de 80\%.}{fig:c03-poder}

\section{O retorno médio é diferente de zero?}
\label{sec:c03-media}

A pergunta mais básica que se pode fazer a um ativo, ou a uma estratégia, é se o retorno esperado é positivo. A média amostral $\bar r$ é uma estimativa dele, e, como toda estimativa, tem erro.

\begin{resultado}{Erro-padrão da média e estatística $t$}
Se $r_1, \dots, r_n$ são não correlacionados, com média $\mu$ e variância $\sigma^2$, então $\Var(\bar r) = \sigma^2/n$. O desvio-padrão de $\bar r$, estimado por $s/\sqrt{n}$, é o \termo{erro-padrão da média}. Para testar $H_0\colon \mu = 0$ usa-se
\begin{equation}
  t = \frac{\bar r}{s/\sqrt{n}} ,
  \label{eq:c03-t-stat}
\end{equation}
que, sob $H_0$, com $n$ grande e dependência fraca entre as observações (o caso dos retornos), tem distribuição aproximadamente normal padrão. Rejeita-se $H_0$ a 5\% quando $|t| > 1,96$.
\end{resultado}

Se os retornos fossem normais e independentes, a distribuição exata de $t$\index{teste t@teste $t$} sob $H_0$ seria a $t$ de Student com $n-1$ graus de liberdade, que é de onde vem o nome da distribuição.\nota{``Student'' era o pseudônimo de William Gosset, que publicou o resultado em 1908.} Com milhares de observações a distinção não tem efeito prático.

\lstinputlisting[style=py, caption={Teste $t$ para a média do log-retorno do Ibovespa (\arq{cap03/l04\_teste\_media.py})}, label=cod:c03-media, firstline=7]{scripts/cap03/l04_teste_media.py}
\begin{lstlisting}[style=saida]
n = 3963 (15.7 anos)   média = 0.0210% ao dia = 5.30% ao ano
desvio = 1.4626%   erro-padrão = 0.0232%
vol. anual = 23.2%   média/vol. anual = 0.228
t = 0.905   p-valor = 0.366
scipy: t = 0.905   p-valor = 0.366
IC de 95% para a média anual: -6.2% a 16.8%
t = 1,96     :  18598 dias =   74 anos
poder de 80% :  38000 dias =  151 anos
\end{lstlisting}

A anualização segue a convenção de 252 dias (Seção~\ref{sec:anualizar}), pela qual os dezesseis anos de calendário da amostra contam como $15,7$. Com todo esse histórico, não se rejeita a hipótese de que o log-retorno esperado do Ibovespa seja zero. O intervalo de confiança\index{intervalo de confianca@intervalo de confiança} de 95\% para a média anual vai de $-6,2\%$ a $16,8\%$: os dados são compatíveis com um índice que perde dinheiro e com um que rende muito bem. As duas últimas linhas fazem a conta inversa. Se a média verdadeira fosse exatamente a amostral, seriam necessários 74 anos de dados para a estatística esperada chegar a $1,96$, e 151 anos para o teste detectar o efeito em 80\% das vezes.

\begin{intuicao}
A estatística $t$ pode ser reescrita em unidades anuais. Com $\mu_a = 252\,\bar r$ e $\sigma_a = s\sqrt{252}$,
\[
  t = \frac{\bar r}{s/\sqrt{n}} = \frac{\mu_a}{\sigma_a}\,\sqrt{\frac{n}{252}} = \frac{\mu_a}{\sigma_a}\,\sqrt{\text{anos}} .
\]
A precisão da média depende do número de anos cobertos pela amostra, qualquer que seja a frequência das observações. Trocar dados diários por dados de cinco minutos multiplica $n$ sem acrescentar um único ano, e o erro-padrão da média anual, $\sigma_a/\sqrt{\text{anos}}$, continua igual.
\end{intuicao}

No Ibovespa, $\mu_a/\sigma_a = 0,228$ e $\sqrt{15,7} = 3,97$, o que dá $t = 0,905$. A razão entre retorno e volatilidade anuais reaparece no Capítulo~\ref{cap:portfolio}, descontada a taxa livre de risco, com o nome de índice de Sharpe. O painel direito da Figura~\ref{fig:c03-poder} generaliza a conta: para uma razão de $0,5$, respeitável para um ativo, dez anos de dados dão poder de 35\%. É por isso que médias históricas de retorno são estimativas tão frágeis, um problema que volta na otimização de carteiras (Capítulo~\ref{cap:portfolio}) e na avaliação de \ingles{backtests} (Capítulo~\ref{cap:backtest}). A volatilidade, em contraste, é estimada com precisão muito maior a partir da mesma amostra.

\section{Testes de normalidade e de autocorrelação}
\label{sec:c03-jb-lb}

\subsection{Jarque-Bera}

O teste de \citet{jarque1987}\index{teste de Jarque-Bera} transforma a comparação visual da Seção~\ref{sec:c03-normal} em uma estatística. A hipótese nula é que os dados vêm de uma distribuição normal, caso em que a assimetria amostral $\hat S$ e o excesso de curtose amostral $\hat K - 3$ devem estar perto de zero.

\begin{resultado}{Estatística de Jarque-Bera}
\begin{equation}
  JB = \frac{n}{6}\left[\hat S^2 + \frac{(\hat K - 3)^2}{4}\right] .
  \label{eq:c03-jb}
\end{equation}
Sob $H_0$ (normalidade) e para $n$ grande, $JB$ tem distribuição qui-quadrado com 2 graus de liberdade. Rejeita-se a normalidade a 5\% quando $JB > 5,99$.
\end{resultado}

A fórmula vem de que, em amostras normais grandes, $\hat S$ tem variância aproximada $6/n$ e $\hat K$, $24/n$. Cada parcela de $JB$ é o quadrado de uma dessas estimativas dividido pela sua variância, e a soma de dois quadrados de normais padrão independentes é uma qui-quadrado com 2 graus de liberdade.

\lstinputlisting[style=py, caption={Jarque-Bera e Shapiro-Wilk em três frequências (\arq{cap03/l05\_normalidade.py})}, label=cod:c03-jb, firstline=10, lastline=25]{scripts/cap03/l05_normalidade.py}
\begin{lstlisting}[style=saida]
diário  n = 3963   JB =  24486.9  p =        0   Shapiro-Wilk: p = 5.1e-40
mensal  n =  192   JB =    241.9  p =    3e-53   Shapiro-Wilk: p = 1e-06
anual   n =   16   JB =      1.2  p =     0.54   Shapiro-Wilk: p = 0.3
scipy, diário: JB = 24486.9
\end{lstlisting}

Para os retornos diários, $JB = 24\,487$ contra um valor crítico de $5,99$; o p-valor é tão pequeno que o computador o arredonda para zero. A normalidade também é rejeitada na frequência mensal, embora, como se viu na Seção~\ref{sec:c03-momentos}, essa rejeição dependa de um único mês. Na frequência anual o teste não rejeita. Leia esse resultado com a armadilha do p-valor em mente: com 16 observações o teste quase não tem poder, a aproximação qui-quadrado é ruim (o teste rejeita menos do que os 5\% nominais), e não rejeitar está longe de demonstrar que retornos anuais são normais.

Outros dois testes aparecem com frequência. O de Shapiro-Wilk\index{teste de Shapiro-Wilk}, na segunda coluna de p-valores, costuma ter bom poder em amostras pequenas; a documentação do SciPy avisa que seu p-valor pode não ser preciso acima de 5\,000 observações. O de Kolmogorov-Smirnov\index{teste de Kolmogorov-Smirnov} compara a função de distribuição empírica com a teórica, mas na forma usual exige que média e variância sejam conhecidas de antemão. Estimá-las da própria amostra, como quase sempre se faz, invalida o p-valor tabelado; nesse caso use a variante de Lilliefors, disponível no \cod{statsmodels}.

\subsection{Ljung-Box}

Para testar se uma série tem autocorrelação, olhar uma defasagem de cada vez leva ao problema da Seção~\ref{sec:c03-multiplos}. O teste de \citet{ljung1978}\index{teste de Ljung-Box} junta as $m$ primeiras autocorrelações amostrais $\hat\rho_1, \dots, \hat\rho_m$ em uma só estatística.

\begin{resultado}{Estatística de Ljung-Box}
\begin{equation}
  Q(m) = n(n+2) \sum_{k=1}^{m} \frac{\hat\rho_k^2}{n-k} .
  \label{eq:c03-lb}
\end{equation}
Sob $H_0$ (observações independentes e identicamente distribuídas), $Q(m)$ tem distribuição aproximada qui-quadrado com $m$ graus de liberdade.
\end{resultado}

\lstinputlisting[style=py, caption={Ljung-Box nos retornos e nos retornos ao quadrado (\arq{cap03/l06\_ljungbox.py})}, label=cod:c03-lb, firstline=11, lastline=15]{scripts/cap03/l06_ljungbox.py}
\begin{lstlisting}[style=saida]
retornos: autocorrelações 1 a 3 = [-0.071  0.033 -0.004]
    lb_stat  lb_pvalue
5   30.0408        0.0
10  68.1284        0.0
20  76.0275        0.0
retornos ao quadrado: autocorrelações 1 a 3 = [0.505 0.527 0.405]
      lb_stat  lb_pvalue
5   3592.4921        0.0
10  4345.8839        0.0
20  4681.2280        0.0
\end{lstlisting}

O resultado dos quadrados é o esperado: $Q(10) = 4\,346$, uma rejeição esmagadora, que confirma o agrupamento de volatilidade. Aplicado a quadrados de retornos ou de resíduos, o teste é o diagnóstico usual desse efeito e dos modelos de volatilidade do Capítulo~\ref{cap:series}.

O resultado dos retornos parece contradizer o segundo fato estilizado, pois o teste também rejeita, com $Q(10) = 68,1$ e autocorrelação de ordem 1 igual a $-0,071$. Duas verificações mostram que a conclusão é mais fraca do que parece. Retirando 2020 da amostra, $Q(10)$ cai para $8,4$, com p-valor de $0,59$: a rejeição vem das semanas de março de 2020, em que quedas de mais de 10\% e repiques de 9\% a 13\% se sucederam em dias próximos. Além disso, a hipótese nula do teste é de observações independentes, e não apenas não correlacionadas. Com volatilidade agrupada, as autocorrelações amostrais variam mais do que $1/\sqrt{n}$, e o teste rejeita com frequência maior do que $\alpha$. Refazendo o teste da autocorrelação de ordem 1 com um erro-padrão que admite variância mutável (o de White, opção \cod{cov\_type="HC0"} do \cod{statsmodels}), a estatística $t$ passa de $-4,49$ para $-1,66$, que não é significativa a 5\%.

\subsection{Uma hipótese por trás de todas as outras}

Todas as contas deste capítulo tratam a amostra como repetições de um mesmo experimento: supõem que média, variância e autocorrelações não mudam com o tempo. Uma série com essa propriedade é dita \termo{estacionária}. Preços não são estacionários, e essa é a razão formal para trabalhar com retornos. O teste de \citet{dickey1979}, em sua versão aumentada (ADF), tem como hipótese nula a presença de uma raiz unitária, isto é, de não estacionariedade. O Capítulo~\ref{cap:series} o apresenta, junto com o que fazer quando a variância dos retornos muda no tempo.

\section{Muitos testes, alguns ``achados''}
\label{sec:c03-multiplos}

Um nível de significância de 5\% significa que, em um teste com $H_0$ verdadeira, há 5\% de chance de alarme falso. Em $m$ testes independentes, todos com $H_0$ verdadeira, a probabilidade de ao menos um alarme falso é $1 - 0,95^m$: 64\% para $m = 20$ e praticamente 100\% para $m = 200$. Quem testa muitas ideias\index{multiplos testes@múltiplos testes} encontra resultados ``significativos'' mesmo que nenhuma funcione.

\figura[t]{cap03/multiplos}{Estatísticas $t$ de 200 séries simuladas de puro ruído, cada uma com cinco anos de retornos diários. Dez delas, em roxo, passam do limite de 5\% sem que haja nada a descobrir. Nenhuma passa do limite de Bonferroni.}{fig:c03-multiplos}

A Figura~\ref{fig:c03-multiplos} mostra uma simulação com 200 ``estratégias'' cujos retornos são números aleatórios de média zero (\arq{cap03/l07\_multiplos.py}). Dez têm $|t| > 1,96$, perto dos 5\% esperados (o número exato varia com a semente). A melhor tem $t = 3,44$, o que em cinco anos corresponde a uma razão anual entre retorno e volatilidade de $1,54$, um número que qualquer gestor apresentaria com orgulho, obtido aqui de um sorteio.

A correção mais simples é a de \termo{Bonferroni}: para manter em $\alpha$ a probabilidade de qualquer alarme falso em $m$ testes, exige-se de cada um p-valor abaixo de $\alpha/m$. Com $m = 200$, o limite sobe de $|t| > 1,96$ para $|t| > 3,66$, e nenhuma das séries simuladas o ultrapassa. O custo é a perda de poder: efeitos verdadeiros e modestos também deixam de ser detectados.

O problema é sério na pesquisa em finanças porque os testes já feitos por outras pessoas, sobre os mesmos dados, também contam. \citet{harvey2016} catalogaram centenas de fatores publicados como explicativos de retornos de ações e argumentam que, dado esse volume de tentativas, um fator novo deveria superar uma estatística $t$ de $3,0$, e não o $2,0$ habitual. O Capítulo~\ref{cap:backtest} retoma o assunto no contexto em que ele mais prejudica: a escolha da melhor entre muitas variantes de uma estratégia.

\section{Covariância e correlação}
\label{sec:c03-corr}

Até aqui cada ativo foi tratado isoladamente. Uma carteira depende de como os ativos se movem em conjunto.

\begin{definicao}{Covariância e correlação de Pearson}
Para duas variáveis $X$ e $Y$ com médias $\mu_X$, $\mu_Y$ e desvios-padrão $\sigma_X$, $\sigma_Y$,
\begin{equation}
  \begin{aligned}
  \Cov(X, Y) &= \E\big[(X-\mu_X)(Y-\mu_Y)\big] , \\
  \rho_{XY} = \Corr(X, Y) &= \frac{\Cov(X, Y)}{\sigma_X\,\sigma_Y} .
  \end{aligned}
  \label{eq:c03-corr}
\end{equation}
\end{definicao}

A covariância é positiva quando desvios do mesmo sinal tendem a ocorrer juntos. Sua unidade é o produto das unidades das duas variáveis, o que a torna difícil de interpretar. A \termo{correlação} divide pelos desvios-padrão e produz um número puro entre $-1$ e $+1$, que não muda se as variáveis forem reescaladas. Os extremos só são atingidos quando uma variável é função linear exata da outra.

O que a correlação de Pearson mede é a intensidade da relação \emph{linear}. Duas consequências são fonte permanente de erro. Variáveis independentes têm correlação zero, mas a recíproca é falsa: correlação zero não implica independência. Os próprios retornos dão o exemplo, pois $r_t$ e $r_{t-1}$ são quase não correlacionados e claramente dependentes, já que o tamanho de um ajuda a prever o tamanho do outro. E a correlação, sendo construída com médias de produtos, é tão sensível a observações extremas quanto a variância.

\figuralarga[t]{cap03/pearson}{Três medidas de associação em quatro situações simuladas. Na relação monótona mas curva, Pearson fica abaixo de 1 e as medidas de posto valem 1. Na dependência em U, em que $y$ é quase determinado por $x$, as três ficam perto de zero. No último painel, 149 pontos sem relação e um único ponto extremo bastam para um Pearson de $0,42$.}{fig:c03-pearson}

Duas alternativas trabalham com a ordem dos dados em vez dos valores. A \termo{correlação de Spearman} é a correlação de Pearson calculada sobre os postos (a posição de cada observação na amostra ordenada). O \termo{tau de Kendall} toma todos os pares de observações e calcula a proporção de pares concordantes, em que $x$ e $y$ variam no mesmo sentido, menos a proporção de discordantes. Ambas medem associação \emph{monótona}, não mudam sob transformações crescentes das variáveis e são pouco afetadas por valores extremos. Como mostra a Figura~\ref{fig:c03-pearson}, nenhuma das três enxerga uma dependência que não seja monótona.

No Código~\ref{cod:c03-corr}, \cod{r} é a tabela de retornos simples diários de quatro ações, de 2015 a 2024. Os métodos \cod{cov} e \cod{corr} de um \cod{DataFrame} devolvem as matrizes de que trata a próxima subseção.

\lstinputlisting[style=py, caption={Covariância, correlação e variância de uma carteira (\arq{cap03/l08\_correlacao.py})}, label=cod:c03-corr, firstline=11]{scripts/cap03/l08_correlacao.py}
\begin{lstlisting}[style=saida]
covariância anualizada:
       PETR4  VALE3  ITUB4  WEGE3
PETR4  0.224  0.092  0.080  0.047
VALE3  0.092  0.191  0.045  0.031
ITUB4  0.080  0.045  0.093  0.034
WEGE3  0.047  0.031  0.034  0.112
correlação de Pearson:
       PETR4  VALE3  ITUB4  WEGE3
PETR4   1.00   0.44   0.55   0.29
VALE3   0.44   1.00   0.34   0.21
ITUB4   0.55   0.34   1.00   0.34
WEGE3   0.29   0.21   0.34   1.00
True
PETR4 x VALE3, pearson : 0.443
PETR4 x VALE3, spearman: 0.385
PETR4 x VALE3, kendall : 0.268
vol. da carteira: 28.3%   média das vols.: 38.8%
\end{lstlisting}

Para PETR4 e VALE3 as três medidas dão $0,44$, $0,39$ e $0,27$. O tau de Kendall é menor em valor absoluto por construção, e não por indicar associação mais fraca: em uma normal bivariada com correlação $\rho$ ele vale $(2/\pi)\arcsin\rho$, que para $\rho = 0,443$ dá $0,29$ \citep{mcneil2015}. Compare cada medida apenas com ela mesma.

\subsection{As matrizes de covariância e de correlação}

Para $N$ ativos com vetor de retornos $\vetor{R}$ e médias $\symbf{\mu}$, as covariâncias de todos os pares são reunidas na \termo{matriz de covariância}
\begin{equation}
  \symbf{\Sigma} = \E\big[(\vetor{R}-\symbf{\mu})(\vetor{R}-\symbf{\mu})\T\big] ,
  \label{eq:c03-sigma}
\end{equation}
simétrica, com as variâncias na diagonal. A \termo{matriz de correlação} tem a mesma estrutura, com 1 na diagonal. A importância de $\symbf{\Sigma}$ está em que ela é tudo o que se precisa para a variância de qualquer carteira. Com pesos $\vetor{w}$, o retorno da carteira é $\vetor{w}\T\vetor{R}$ (Equação~\eqref{eq:ret-carteira}) e
\begin{equation}
  \Var(\vetor{w}\T\vetor{R}) = \vetor{w}\T\symbf{\Sigma}\,\vetor{w} = \sum_{i=1}^{N}\sum_{j=1}^{N} w_i w_j \sigma_i \sigma_j \rho_{ij} .
  \label{eq:c03-var-cart}
\end{equation}
Como uma variância não pode ser negativa, $\vetor{w}\T\symbf{\Sigma}\,\vetor{w} \ge 0$ para todo $\vetor{w}$: a matriz é positiva semidefinida. A última linha da saída do Código~\ref{cod:c03-corr} aplica a fórmula a quatro ações com pesos iguais. A volatilidade da carteira, de $28,3\%$ ao ano, é bem menor do que a média das volatilidades individuais, de $38,8\%$, porque as correlações estão longe de 1. Essa é a aritmética da diversificação, e o Capítulo~\ref{cap:portfolio} parte dela para chegar à fronteira eficiente de \citet{markowitz1952}.

\figura[t]{cap03/matriz_corr}{Correlações de Pearson entre retornos diários de sete ações brasileiras e do S\&P~500, de 2015 a 2024, nos $2\,409$ dias com pregão nos dois países. Todas são positivas. Os dois bancos formam o par mais correlacionado ($0,76$), e as correlações entre ações vão de $0,21$ a $0,76$, com média de $0,36$. Alguns valores diferem na segunda casa dos do Código~\ref{cod:c03-corr}, cuja amostra inclui os dias sem pregão em Nova York.}{fig:c03-matriz}

Uma matriz com $N$ ativos tem $N(N+1)/2$ elementos distintos: 28 para as sete ações da Figura~\ref{fig:c03-matriz}, $5\,050$ para uma carteira de cem. Cada um é estimado com erro, a partir de uma amostra que não cresce com $N$, e a matriz amostral se torna pouco confiável quando o número de ativos se aproxima do número de observações. Um remédio usual é o encolhimento (\ingles{shrinkage}) proposto por \citet{ledoit2004}, e o problema volta no Capítulo~\ref{cap:portfolio}.

\begin{armadilha}{Correlacionar séries dessincronizadas}
De 2015 a 2024, no arquivo \arq{precos\_ajustados.csv}, a correlação entre o retorno diário do Ibovespa e o do dólar (\cod{BRL=X}) na mesma data é $0,01$. Deslocando o dólar em um dia, ela vai a $-0,46$ (\arq{cap03/c03\_sincronia.py}). A relação sempre esteve lá: nessa série, a variação do dólar registrada em uma data corresponde ao que a B3 viu no pregão anterior. A Seção~\ref{sec:c02-cambio} documenta o defeito, com números um pouco diferentes ($-0,02$ e $-0,41$) porque usa o período de 2010 a 2025, e mostra que a PTAX do Banco Central não o tem. Antes de interpretar uma correlação entre mercados diferentes, confira como cada série marca a data e a hora do fechamento, e calcule também as correlações defasadas.
\end{armadilha}

\section{Correlação que muda no tempo}
\label{sec:c03-movel}

Um único número para dez anos esconde variação. A \termo{correlação móvel} recalcula a correlação a cada dia usando apenas as últimas $J$ observações; em pandas, \cod{r["A"].rolling(63).corr(r["B"])}. A Figura~\ref{fig:c03-movel} faz isso com janela de 63 pregões (um trimestre) para os 21 pares formados pelas sete ações e mostra a média, ao lado da volatilidade do Ibovespa na mesma janela.

\figura[t]{cap03/corr_movel}{Volatilidade do Ibovespa e correlação média entre sete ações, em janelas móveis de 63 pregões. A correlação média oscila entre $0,09$ e $0,78$, com mediana de $0,32$. O máximo, em 18 de março de 2020, coincide com o salto da volatilidade.}{fig:c03-movel}

Entre 20 de fevereiro e 30 de abril de 2020, a correlação média entre as sete ações foi de $0,71$, contra $0,36$ no período inteiro. A de PETR4 com VALE3, que em dez anos é $0,44$, foi de $0,83$ nesse intervalo; a correlação móvel de 63 pregões do par chegou a $0,89$ em 16 de março. Para quem monta carteiras a consequência é desagradável: a diversificação medida em tempos normais superestima a proteção disponível justamente quando ela é mais necessária.

Três cuidados impedem que essa observação vire uma conclusão maior do que os dados sustentam.

O primeiro é o erro de estimação. Para dados normais, o erro-padrão de uma correlação amostral é aproximadamente $(1-\rho^2)/\sqrt{n}$. Com $\rho = 0,3$ e $n = 63$ isso dá $0,11$, ou uma faixa de $\pm 0,22$ em torno do valor verdadeiro.\nota{\arq{cap03/c03\_fator.py} confere por simulação essa aproximação e o exemplo do fator comum.} Boa parte do sobe e desce de uma correlação móvel entre dois ativos é ruído de amostra, e a média de 21 pares da figura é mais estável do que qualquer par.

O segundo é que a amostra contém uma única crise dessa magnitude. Ela ilustra o fenômeno, mas não permite medi-lo com precisão. A evidência de que correlações entre mercados de ações aumentam nas quedas acentuadas vem de estudos com amostras muito mais longas, como o de \citet{longin2001}, com retornos mensais de 1959 a 1996 dos cinco maiores mercados. Eles encontram correlação maior nas grandes quedas, mas não nas grandes altas.

O terceiro é mais sutil: a correlação pode subir sem que a ligação entre os ativos tenha mudado.

\begin{exemplo}{Correlação e volatilidade de um fator comum}
Suponha que duas ações sigam $R_i = M + \varepsilon_i$, em que $M$ é um movimento de mercado com desvio-padrão $\sigma_M$ e $\varepsilon_1$, $\varepsilon_2$ são componentes específicos, independentes entre si e de $M$, com desvio-padrão $\sigma_\varepsilon$. Então $\Cov(R_1, R_2) = \sigma_M^2$, $\Var(R_i) = \sigma_M^2 + \sigma_\varepsilon^2$ e
\[
  \rho = \frac{\sigma_M^2}{\sigma_M^2 + \sigma_\varepsilon^2} .
\]
Com $\sigma_\varepsilon = 1,5\%$ e $\sigma_M = 1\%$ ao dia, $\rho = 1/3,25 = 0,31$. Se a volatilidade do mercado quadruplica, para $\sigma_M = 4\%$, e nada mais muda, $\rho = 16/18,25 = 0,88$.
\end{exemplo}

Nesse modelo a sensibilidade de cada ação ao mercado é a mesma nos dois regimes; o que muda é a proporção da variância que vem do fator comum. \citet{forbes2002} mostraram que esse efeito enviesa para cima as correlações medidas em períodos turbulentos e que, depois de corrigi-lo, boa parte do aparente ``contágio'' entre mercados em crises dos anos 1980 e 1990 desaparece. Uma correlação alta em 2020 não exige, por si só, um mecanismo novo de dependência. O exemplo, porém, é simétrico: a volatilidade do fator sobe tanto nas quedas quanto nas altas, e sozinho ele não explica a diferença entre umas e outras encontrada por \citet{longin2001}.

\begin{repositorio}
\arq{aula3/s3\_03\_testes\_hipotese.py} aplica o teste $t$, o de Ljung-Box e os de normalidade, além de uma comparação de duas amostras. \arq{aula3/s3\_04\_correlacao.py} calcula as três correlações, a correlação móvel e a comparação entre períodos de crise e de calma para um conjunto maior de ativos.
\end{repositorio}

\begin{resumo}
  \item A média diária de um retorno é dezenas de vezes menor do que seu desvio-padrão. Por isso o retorno esperado é estimado com enorme incerteza, que só diminui com mais anos de dados, e não com dados mais frequentes.
  \item Retornos diários têm caudas pesadas, autocorrelação linear quase nula e volatilidade agrupada. A normal subestima a frequência de dias extremos por ordens de grandeza; a $t$ de Student com poucos graus de liberdade descreve melhor as caudas, mas não o agrupamento.
  \item Assimetria, curtose e os testes construídos sobre elas são dominados por poucas observações. Refaça a conta sem o período extremo antes de acreditar nela.
  \item O p-valor é calculado supondo $H_0$ verdadeira. Ele não é a probabilidade de $H_0$, e não rejeitar não equivale a confirmar, sobretudo quando o poder é baixo.
  \item Testar muitas hipóteses ao nível de 5\% torna os falsos achados quase certos. O limite de significância precisa subir com o número de tentativas.
  \item A correlação de Pearson mede apenas dependência linear, é sensível a extremos e muda no tempo. Ela sobe nas crises, em parte porque a volatilidade do fator comum sobe.
\end{resumo}

\begin{exercicios}
  \item \nivel{básico} Carregue os preços ajustados de \cod{WEGE3.SA} e \cod{MGLU3.SA} de 2015 a 2024 e monte, para os log-retornos diários, a tabela do Código~\ref{cod:c03-descritiva} acrescida de assimetria e excesso de curtose. Em qual dos dois ativos a diferença entre $s$ e $\text{IQR}/1,349$ é maior?
  \item \nivel{básico} Uma $t$ de Student tem $\nu = 6$ graus de liberdade. Calcule a variância e o excesso de curtose. Qual fator de escala a transforma em uma variável de variância 1? Confira com \cod{scipy.stats.t.stats(6, moments="vk")}.
  \item \nivel{básico} Uma estratégia rendeu em média 8\% ao ano, com volatilidade de 16\% ao ano, durante 6 anos. Calcule a estatística $t$ para a hipótese de retorno esperado nulo. Quantos anos seriam necessários para chegar a $t = 1,96$ com a mesma média e a mesma volatilidade?
  \item \nivel{intermediário} Faça o QQ-plot dos log-retornos diários de \cod{PETR4.SA} contra a normal e contra a $t$ ajustada por máxima verossimilhança. Repita com retornos mensais. O que muda, e qual fato estilizado explica a mudança?
  \item \nivel{intermediário} Implemente a estatística da Equação~\eqref{eq:c03-lb} com NumPy, para $m = 10$, e confira com \cod{acorr\_ljungbox} nos retornos e nos retornos ao quadrado de \cod{VALE3.SA}. Repita excluindo 2020 e comente.
  \item \nivel{intermediário} Simule $1\,000$ amostras de 250 observações de uma $t$ com $\nu = 4$ e aplique o teste de Jarque-Bera a cada uma. Em que fração das amostras a normalidade é rejeitada a 5\%? Repita com amostras de 25 observações. Que conceito da Seção~\ref{sec:c03-teste} os dois resultados ilustram?
  \item \nivel{desafio} Refaça a simulação da Figura~\ref{fig:c03-multiplos} trocando o ruído normal por uma $t$ com $\nu = 3$ reescalada para o mesmo desvio-padrão. A proporção de falsos positivos do teste $t$ a 5\% muda? Relacione o resultado com o teorema central do limite.
  \item \nivel{desafio} Verifique por simulação o exemplo do fator comum: gere $R_1$ e $R_2$ com $\sigma_M$ alternando entre 1\% e 4\% em blocos de 250 dias e $\sigma_\varepsilon = 1,5\%$. Calcule a correlação em cada regime e a correlação da amostra inteira. Por que esta última não é a média das duas?
\end{exercicios}

\begin{leituras}
  \item \citet{cont2001}, o artigo inteiro, curto e legível, para a lista completa de fatos estilizados e os cuidados estatísticos ao medi-los.
  \item \citet{tsay2010}, capítulo 1, para momentos amostrais, testes de normalidade e distribuições de retornos com exemplos de dados reais.
  \item \citet{campbell1997}, os capítulos iniciais, para os testes de previsibilidade, incluindo versões que admitem heterocedasticidade.
  \item \citet{mcneil2015}, capítulos 3 e 7, para caudas pesadas e para os limites da correlação como medida de dependência.
  \item \citet{harvey2016}, para múltiplos testes na pesquisa em finanças.
\end{leituras}
